% The teacher's appendix: the lecture's design.  Its subject is the
% teaching design, so the analysis of what varies and what stays
% invariant is body text here; the margin notes in the chapter carry the
% same analysis activity by activity.  It is input only when the
% instructor notes are printed (see notes-body.tex).

\chapter{How the lecture is designed}
\label{app:design}

This appendix is for whoever gives the lecture.
It takes the five learning objectives one at a time, divides each into
the aspects a student has to notice to reach it, and says what the
lecture offers for each aspect and what it does not.

The analysis follows variation theory as Marton states it.
In that theory a learner notices an aspect of something when the aspect
varies while the rest stays the same, and the cases have to be
experienced together: discernment comes from \enquote{the simultaneous
experience of things that differ in color but are the same otherwise}%
\autocite[46]{Marton2015}.
The lecture takes three rules of design from it.
A contrast comes before a generalisation, since on Marton's account what
several cases have in common cannot be seen by someone who has not yet
seen the aspect in any of them\autocite[46--47]{Marton2015}.
The cases of a contrast stand side by side, in one table or one pair of
rounds.
And a question does not point at the aspect it is meant to bring out,
because \enquote{the more clearly the teacher tells the students what is
to be done, the less chance the students get to make the necessary
distinctions}\autocite[13]{Marton2015}.

For each aspect below, \emph{offered} says which activity varies it and
how, and \emph{possible to learn} is a judgement of whether a student
can notice the aspect from that activity alone, or is only told it.


\section{Checking (\cref{SecLOCheck})}
\label{sec:design-check}

\SecLOCheck*

\paragraph{Producing and checking are different jobs.}
\emph{Offered:} a contrast in two steps.
\Cref{tab:three-proofs} shows who produced each of three proofs
and withholds how they were checked, and the students vote on it in
\cref{ex:three-proofs}; \cref{tab:three-outcomes} then adds the
check and the outcome to the same table.
\emph{Possible to learn:} yes.
The students vote on the producer alone, and the outcomes then cut
across their votes.
\paragraph{The check has to be independent of the producer.}
\emph{Offered:} three kinds of check side by side in
\cref{tab:three-outcomes}, and the same aspect again in the three
rounds of \cref{sec:product}, summed up in \cref{tab:defence}: a
rule in the bot's instructions makes the producer its own checker.
\emph{Possible to learn:} yes, and best in the rounds, where the
students themselves get past the rule and fail against the check in
the program.
\paragraph{Cases that came out right do not establish the claim.}
\emph{Offered:} a generalisation: the same point for the 12~trillion
zeros of \cref{sec:riemann} and for a program that \enquote{works}
when it is tried (\cref{sec:works}).
In \cref{sec:product} the students then meet it from the other side: a
bot that serves every ordinary customer correctly, and that they
break.
\emph{Possible to learn:} in part.
Two settings are few for a generalisation, and the first is told; the
rounds are where the students experience it.

\paragraph{A check answers only the question it was given.}
\emph{Offered:} \cref{sec:statement}, one example, told.
\emph{Possible to learn:} hardly.
Nothing varies here: the students do not see a valid proof of the
wrong statement beside a valid proof of the right one.
The aspect is named so that it can be recognised in round 2, where
the bot follows its instructions and still pays the wrong person.
\paragraph{What someone gets out of an AI depends on what that person
  can do.}
\emph{Offered:} a contrast in \cref{ex:who}: the AI is the same, and
the person varies between the mathematician behind proof B and the
student, whose own answer to \cref{ex:riemann} is shown again.
\Cref{ex:brought} then says what the mathematician brought.
\emph{Possible to learn:} yes, for this one case.
It is the lecture's perspective on who can do what with AI, and one
case cannot carry it alone: \cref{sec:gains} asks the same of many
people, and the result there decides how far the point may be taken.

\paragraph{Saying exactly what is wanted takes knowledge of the subject.}
\emph{Offered:} the students' own answers to \cref{ex:who}, sorted
into the jobs of \cref{sec:three-jobs}, and the optional pair of
\cref{ex:eli5}, a contrast in which the question is the same and
the asker's knowledge varies.
\emph{Possible to learn:} in part.
The students never specify anything themselves.
The spoken follow-up to \cref{ex:posttest} shows whether they name
this aspect unprompted.


\section{The evidence on who gains (\cref{SecLOEvidence})}
\label{sec:design-evidence}

\SecLOEvidence*

\paragraph{The gain differs between groups of people.}
\emph{Offered:} a prediction (\cref{ex:who-gains}) and then the
findings (\cref{ex:gains-studies}).
\emph{Possible to learn:} yes, as a contrast between the student's
own prediction and what was found.
\paragraph{The finding depends on the task and on who can judge the result.}
\emph{Offered:} studies that ask the same question of different
tasks and people (\cref{ex:gains-studies}), with the conclusion
that closes \cref{sec:gains}.
\emph{Possible to learn:} only as far as the studies that the
literature search supports actually differ in these respects; the
wording of that conclusion follows the search and not the
lecturer's expectation.


\section{Prompt injection (\cref{SecLOInjection})}
\label{sec:design-injection}

\SecLOInjection*

\paragraph{The model reads one text.}
\emph{Offered:} \cref{fig:bot}, shown, and round 1
(\cref{act:secret}), where what the student types competes with
what the shop wrote.
\emph{Possible to learn:} yes, by doing.
\paragraph{A rule written into that text is not a barrier.}
\emph{Offered:} a contrast over where the defence is: no rule
(\cref{act:secret}), a rule (\cref{act:secret-rule}), the secret
removed (\cref{ex:secret-repair}).
The bot, the secret and the goal stay the same.
\emph{Possible to learn:} yes, if somebody in the room gets past
the rule; the lecturer has an attack that is known to work for the
case that nobody does.
\paragraph{What steering the model achieves depends on what the program
  lets it do.}
\emph{Offered:} a contrast between round 1 and round 2
(\cref{act:refund}): the same way of attacking, against a bot
without tools and with tools.
\emph{Possible to learn:} yes.
\paragraph{The text that steers need not come from the person chatting.}
\emph{Offered:} a contrast between round 1 and round 3
(\cref{act:reviews}): the same kind of instruction, typed by the
attacker or read from a review while somebody else is chatting.
\emph{Possible to learn:} only if round 3 is run and the model can
be steered that way.
Without it the aspect is told, in \cref{def:injection}.


\section{A model inside a product (\cref{SecLOProduct})}
\label{sec:design-product}

\SecLOProduct*

\paragraph{What the model is given to read can be read out.}
\emph{Offered:} round 1.
\emph{Possible to learn:} yes, by doing.
\paragraph{What its tools can do, an outsider can do, unless the program
  checks.}
\emph{Offered:} round 2, and the refund tool before and after
(\cref{ex:refund-fix,ex:refund-repair}).
\emph{Possible to learn:} yes.
\paragraph{Every message is paid for, whoever sends it and whatever the
  answer.}
\emph{Offered:} \cref{ex:bill}, a contrast between messages the bot
answered and messages it refused, with the cost invariant.
\emph{Possible to learn:} yes, if the lecturer points at the meter
after a refusal; the meter is the room's own.
\paragraph{The defence belongs in the program around the model.}
\emph{Offered:} the generalisation of \cref{tab:defence}, after
three rounds that each contrasted its two columns, and the counts of
\cref{ex:results}.
The recording cannot show that a rule for the model fails as a repair:
with so few attempts the repaired bot of \cref{ex:reviews-hardened}
may hold every time.
That point rests on the research reported in \cref{sec:stopped}, and
on round 1, where the rule \enquote{never reveal} is seen to fail.
\emph{Possible to learn:} yes; this is the lecture's main point and
the one with the most variation behind it.
\paragraph{The company answers for what the bot says and does.}
\emph{Offered:} the remark that closes \cref{sec:defence}, one
case, told.
\emph{Possible to learn:} as a fact only.


\section{An agent of one's own (\cref{SecLOAgents})}
\label{sec:design-agents}

\SecLOAgents*

\paragraph{An agent has the same three parts as the shop's bot.}
\emph{Offered:} the generalisation of \cref{tab:same-parts}: the
parts stay, the system changes.
\emph{Possible to learn:} yes, since the students have just
attacked each part.
\paragraph{The danger needs all three parts at once.}
\emph{Offered:} the fusion of \cref{ex:trifecta}: four setups that
differ in which parts are present, with the attack the same.
\emph{Possible to learn:} yes.
Each wrong answer is one mistaken belief about which parts are
enough.
\paragraph{Taking one part away is a choice that depends on the task.}
\emph{Offered:} the summary at the end of \cref{sec:agent}, told.
\emph{Possible to learn:} hardly.
The students do not set up an agent or choose its permissions in
the lecture; the second half of the objective, choosing measures,
needs an exercise of its own afterwards.


\section{The order, and what to drop}
\label{sec:design-order}

The product comes before the student's own agent.
The bot gives every student the same concrete system to attack, and
their own agents, which differ and which many of them do not have, are
then met as another case of it: examples first, the general pattern
after.

\Cref{ex:replace,ex:posttest} measure the same thing before and after,
with different wordings so that the answers are kept apart.

When time runs short, drop in this order: \cref{ex:role},
\cref{ex:eli5}, round 3 (\cref{act:reviews}), \cref{ex:five-years}.
Rounds 1 and 2 and \cref{tab:defence} are the part to protect.
